% ============================================================
% Section 5: Conclusion
% BearMamba-3: Mamba-3 + L_kin + Multi-Sensor Fusion
% Target: IEEE TIM (IF 5.6)  Draft: 2026-06-15 → 2026-06-18 (D34)
% ============================================================

\section{Conclusion}
\label{sec:conclusion}

We presented BearMamba-3, a bearing fault diagnosis framework that integrates
the Mamba-3 state-space backbone with a kinematic frequency inductive bias
$\mathcal{L}_{\text{kin}}$ and multi-sensor convolutional fusion.

The core structural insight is that Mamba-3's complex oscillatory states carry
instantaneous rotation frequencies $f_{h,k}(t)$ — a physical quantity absent
in real-valued SSMs — which enable activation-level alignment with analytically
known bearing fault harmonics.
This binding between the backbone selection (C1) and the physical regulariser
design (C2) constitutes a principled, integrated contribution rather than two
independently interchangeable components.

Our experimental results on three public datasets reveal a coherent picture
of where activation-level kinematic priors help and where they do not:
$\mathcal{L}_{\text{kin}}$ consistently reduces prediction variance under
synthetic AWGN (CWRU, std $\downarrow$14--62\%) and aligns internal state
frequencies with fault harmonics in both constant-speed (CWRU) and
run-to-failure (XJTU-SY) regimes.
However, it provides no statistically significant mean accuracy improvement,
and is ineffective in discrete-speed-switching regimes (Paderborn) where the
per-sample RPM degeneracy eliminates its continuous-speed adaptation advantage.
These boundary conditions are reported as honest negative results alongside
the positive findings.

The multi-sensor experiments (C3) reveal a precision--sample-efficiency
trade-off: dual-sensor fusion improves accuracy when training data is
sufficient to learn cross-channel correspondences (CWRU $-4$~dB, $p=0.031$;
XJTU cross-condition IR recall $+31.6$~pp), but induces negative transfer in
leave-one-bearing-out scenarios with a single training instance per fault type.
This finding has practical implications for industrial deployment: multi-sensor
systems require adequate per-fault-type coverage, a constraint that may be
difficult to satisfy in early commissioning or rare-fault scenarios.

\paragraph{Inductive bias paradigms.}
A cross-dataset analysis (Section~\ref{subsec:paradigm}) further reveals a
systematic complementarity between two families of diagnostic methods.
\emph{BN-dependent} methods (1D-CNN family) use BatchNorm to normalise
activations using source-domain statistics, conferring powerful within-condition
regularisation but collapsing under domain shift (XJTU-SY OOD, $41.3\%$
macro-F1, IR recall $< 2\%$).
The \emph{frequency-grounded} paradigm introduced here substitutes statistical
normalisation with an explicit, distribution-agnostic physical constraint,
yielding comparable within-condition accuracy (the $10.2$~pp gap at $-8$~dB
closes to $p=0.125$ once BatchNorm is removed from 1D-CNN) and maintaining
macro-F1 of $79.7\%$ on the XJTU-SY cross-condition benchmark ($+38.4$~pp
over 1D-CNN).
This paradigm distinction is the paper's principal generalisable insight:
the deployment context --- within-condition vs.\ cross-condition monitoring ---
should govern the choice of inductive bias, not accuracy on a single benchmark.

\paragraph{Limitations and future work.}
The MIMO backward limitation on Ampere-class GPUs (sm\_86) prevented the full
validation of the hypothesised MIMO--sensor-count interaction; future work on
Hopper-class hardware could provide a complete ablation.
Non-Gaussian industrial noise remains an open direction, as detailed
in the Discussion (Section~\ref{sec:discussion}).
Self-supervised pretraining~\cite{ding2022ssl,wan2022siamese} could further
reduce label requirements in run-to-failure monitoring scenarios such as
XJTU-SY, complementing the supervised protocol used here.
Larger-scale replication studies (seeds $> 5$) would strengthen the
variance-stability and multi-sensor findings, where current statistical
power is $\approx 50\%$ for 2~pp effects at $n=5$.

Beyond these specific extensions, BearMamba-3 provides a
measurement-stage building block for intelligent diagnostic
instrumentation: it pairs an interpretable, physically grounded
state-space representation with an honest assessment of the regimes
in which physical inductive bias does and does not confer accuracy
gains, supporting evidence-based selection of diagnostic algorithms
in industrial condition-monitoring systems.
